\chapter[Fundamentação Teórica]{Fundamentação Teórica}
\label{cap:fundamentacao-teorica}

Este capítulo apresenta os conceitos necessários para compreender as escolhas metodológicas do trabalho. A intenção não é fazer uma revisão geral de aprendizado profundo, mas explicar os elementos que aparecem diretamente nos modelos avaliados na comparação principal: verificação por similaridade, embeddings faciais, redes siamesas, backbones convolucionais e Transformers, mecanismos de atenção entre faces, adaptação por LoRA, recuperação de exemplos, modelos multimodais e métricas de avaliação.

\section{Formulação do problema}

A verificação visual de parentesco facial pode ser formulada como uma tarefa de classificação binária. Dado um par de imagens de rosto, $x_a$ e $x_b$, o sistema deve decidir se as duas pessoas possuem ou não vínculo de parentesco. Em termos computacionais, um modelo aprende uma função $f_\theta : \mathbb{R}^{H \times W \times 3} \to \mathbb{R}^d$ que transforma cada face em um vetor de características, ou embedding, de dimensão $d$. A partir desses vetores, calcula-se um escore de similaridade e compara-se esse escore com um limiar de decisão $\tau$.

Neste trabalho, quando a comparação usa cosseno, o escore do par é dado por $s(x_a, x_b) = \cos(f_\theta(x_a), f_\theta(x_b))$. A decisão fica concentrada no limiar: acima de $\tau$, o par é aceito como parentesco; abaixo desse valor, é rejeitado como não-parentesco. A dificuldade não está nessa regra final, mas em produzir uma ordenação confiável antes dela. Parentes podem aparecer com idades, expressões, iluminação e poses bastante diferentes, enquanto pessoas sem vínculo biológico às vezes têm rostos parecidos.

Os modelos supervisionados deste trabalho resolvem essa versão binária do problema. No estudo consolidado com modelos multimodais visão-linguagem, ou VLMs, a pergunta também é binária: decidir entre kin e non-kin em regime zero-shot. Uma avaliação anterior com VLMs tentou escolher uma entre as onze relações positivas do FIW, como pai-filho, mãe-filha, irmãos ou avô-neto, mas essa tarefa multiclasse é tratada apenas como diagnóstico histórico, porque não inclui a classe non-kin e mede uma pergunta diferente.

\section{Redes siamesas e aprendizado por similaridade}

Redes siamesas são uma forma natural de modelar tarefas baseadas em pares. A arquitetura tem dois ramos idênticos, com pesos compartilhados, e cada ramo processa uma imagem do par. Como os pesos são compartilhados, ambas as imagens passam pela mesma transformação e são mapeadas para um espaço vetorial comum. Em seguida, os vetores resultantes são comparados por uma medida de similaridade ou distância.

A ideia por trás desse tipo de rede é aproximar pares positivos e afastar pares negativos no espaço vetorial. Em verificação facial, essa estrutura é útil porque a pergunta não é qual é a classe isolada de uma imagem, mas se duas imagens devem ser consideradas próximas segundo algum critério. No caso do parentesco, o critério não é identidade da pessoa, e sim semelhança familiar. Essa adaptação para verificação de parentesco foi consolidada no trabalho de \citeonline{robinson2018visual}, que usou redes siamesas com backbones pré-treinados em reconhecimento facial.

A similaridade de cosseno é utilizada porque avalia o alinhamento angular entre embeddings, independentemente da magnitude, ou norma $L^2$, dos vetores. Com normalização em $L^2$, seu valor fica no intervalo $[-1, 1]$. Nessa condição, a distância euclidiana e o cosseno dão leituras próximas do mesmo espaço, pois ambas dependem da posição relativa entre os vetores. Já a margem aparece nas perdas discriminativas como a distância mínima que se deseja manter entre positivos e negativos.

\section{Redes convolucionais, ResNet e ConvNeXt}

Redes neurais convolucionais, ou CNNs, foram durante muitos anos a base principal dos modelos de visão computacional. Uma CNN aprende filtros locais que percorrem a imagem de entrada e produzem mapas de características. Nas primeiras camadas, esses filtros tendem a capturar padrões simples, como bordas e texturas. Nas camadas mais profundas, as ativações passam a representar partes mais complexas da imagem, como regiões da face e estruturas semânticas.

Em reconhecimento facial, arquiteturas residuais tiveram papel central. A ResNet \cite{he2016deep}
introduziu conexões de atalho que facilitam o treinamento de redes profundas, reduzindo o problema de degradação do gradiente em muitas camadas. Essa família de arquitetura é importante para este trabalho porque aparece tanto em modelos clássicos de reconhecimento facial quanto em backbones especializados usados como extratores de embeddings.

O AdaFace \cite{kim2022adaface} é um exemplo de backbone facial especializado. Ele usa uma variante IR-101 da família ResNet adaptada para faces e é treinado em larga escala, no WebFace4M, para separar identidades. A perda do AdaFace ajusta a margem de acordo com a qualidade da imagem, dando tratamento diferente a exemplos mais confiáveis e a exemplos degradados. O resultado é um extrator que produz embeddings de 512 dimensões em que distâncias e ângulos carregam informação de similaridade facial. Neste trabalho, esse sinal é usado de duas formas: diretamente na linha de base AdaFace-Cosseno e como base do AdaFace-Regional, com descongelamento parcial do último estágio.

ConvNeXt \cite{liu2022convnet} revisita a CNN tradicional incorporando escolhas de projeto associadas aos Transformers, como convoluções depthwise, blocos invertidos e normalização por camada. No escopo deste trabalho, ele aparece apenas no experimento complementar ConvNeXt-ViT.

\section{Vision Transformers e DINOv2}

O Vision Transformer \cite{dosovitskiy2020image} aplica a arquitetura Transformer, originalmente proposta para processamento de texto \cite{vaswani2017attention},
diretamente sobre imagens. Em vez de usar filtros convolucionais, a imagem é dividida em pequenos patches não-sobrepostos. Cada patch é projetado para um vetor, recebe informação de posição e passa a compor uma sequência de tokens processada pelo encoder Transformer.

O mecanismo central do Transformer é a self-attention, definida por

\begin{equation}
\text{Attention}(Q, K, V) = \mathrm{softmax}\!\left( \frac{QK^\top}{\sqrt{d_k}} \right) V,
\label{eq:attention}
\end{equation}

onde $Q$, $K$ e $V$ são projeções lineares dos tokens de entrada, chamadas respectivamente de query, key e value, e $d_k$ é a dimensão de cada cabeça de atenção. A atenção permite que cada token compare sua informação com todos os outros tokens da sequência. Em imagens faciais, isso é relevante porque regiões distantes, como olhos, boca e contorno da mandíbula, podem ser relacionadas diretamente sem depender apenas do crescimento gradual do campo receptivo.

DINOv2 \cite{oquab2024dinov2}
é um Vision Transformer treinado de forma auto-supervisionada em grande escala, sobre 142 milhões de imagens sem rótulos. O método combina destilação entre diferentes recortes aumentados de uma mesma imagem e objetivo de máscara em nível de patch, produzindo representações gerais que podem ser transferidas para outras tarefas. Para parentesco facial, a hipótese de interesse é que uma representação visual ampla pode conter informação útil sobre estrutura facial, mesmo sem ter sido treinada especificamente para parentesco.

No DINOv2-LoRA deste trabalho, o DINOv2 funciona como backbone de uso geral e é adaptado ao domínio por módulos treináveis e por diferentes regimes de descongelamento, conforme descrito no Capítulo~\ref{cap:Metodologia}. Essa escolha contrasta com o AdaFace, que já parte de um pré-treino especializado em reconhecimento facial. A comparação entre esses caminhos ajuda a separar dois tipos de conhecimento visual: representações gerais aprendidas em larga escala e representações faciais treinadas para discriminar identidades.

\section{Atenção cruzada entre faces}

Atenção cruzada, ou cross-attention, é uma variação da self-attention em que as consultas vêm de uma sequência e as chaves e valores vêm de outra. Em verificação de parentesco, isso permite calcular pesos de atenção entre os tokens de uma face e os tokens da outra. Em vez de comparar apenas dois embeddings globais, o módulo estima relações entre partes de uma imagem e partes da outra.

Formalmente, dados $X_a$ e $X_b$ como os tokens das duas faces, a atenção cruzada de $a$ sobre $b$ pode ser escrita como

\begin{equation}
\text{CrossAttn}(X_a, X_b) = \mathrm{softmax}\!\left( \frac{X_a W_Q (X_b W_K)^\top}{\sqrt{d_k}} \right) X_b W_V,
\label{eq:cross_attention}
\end{equation}

onde $W_Q$, $W_K$ e $W_V$ são matrizes aprendíveis. A versão bidirecional calcula esse tipo de interação nos dois sentidos, de $a$ para $b$ e de $b$ para $a$. O efeito prático é permitir que uma face influencie a representação da outra durante a comparação. Esse é o princípio central do FaCoR \cite{su2023facornet} e da implementação ViT-FaCoR avaliada neste trabalho.

Uma limitação dos modelos baseados em patches densos é que a atenção pode envolver muitas posições pouco informativas para parentesco, como fundo, cabelo ou regiões afetadas por iluminação. Por isso, este trabalho também explora uma formulação com regiões anatômicas fixas. Em vez de usar todos os patches da imagem, a face alinhada é dividida em um conjunto pequeno de regiões, como face global, olhos, nariz, boca e mandíbula. Cada região é transformada em um token, e a atenção cruzada passa a operar sobre uma matriz pequena, de cinco por cinco no caso do AdaFace-Regional. A formulação completa combina esses tokens com um gate sigmoide aprendido por região, permitindo que mais de uma parte da face seja considerada relevante ao mesmo tempo.

Essa escolha introduz um prior estrutural simples: o modelo é orientado a comparar partes faciais interpretáveis, e não todos os pontos da imagem com a mesma liberdade. A vantagem é reduzir a complexidade da comparação e aproximar a arquitetura do modo como o problema é descrito visualmente. O custo é uma dependência maior do alinhamento facial. Se a face é mal detectada, se a pose é extrema ou se o recorte desloca as regiões, os tokens deixam de corresponder às partes anatômicas esperadas. Na revisão feita para este trabalho, não foi identificada uma combinação equivalente de backbone facial, tokens anatômicos, atenção cruzada regional e passagem simétrica aplicada à verificação de parentesco facial.

Há ainda um detalhe específico da tarefa binária: inverter a ordem das imagens não deveria mudar o rótulo. Se ``A é parente de B'' é verdadeiro, então ``B é parente de A'' também deve ser. Em termos ideais, a função de decisão deveria satisfazer $f(A,B) = f(B,A)$. Essa propriedade, porém, não aparece automaticamente quando a cabeça recebe embeddings concatenados em ordem fixa ou quando a atenção cruzada é usada de forma direcional. A passagem simétrica processa o mesmo par nas duas ordens e regulariza a função de decisão, reduzindo a chance de dependência de atalhos associados à posição da imagem no par.

Uma variação recente do mecanismo de atenção é a atenção diferencial \cite{ye2024differential},
que divide cada cabeça em duas partes e calcula a atenção como a diferença entre dois mapas softmax:

\begin{equation}
\text{DiffAttn}(Q, K, V) = \big( \mathrm{softmax}(Q_1 K_1^\top / \sqrt{d}) - \lambda \, \mathrm{softmax}(Q_2 K_2^\top / \sqrt{d}) \big) V,
\label{eq:diff_attn}
\end{equation}

onde $\lambda$ é um escalar aprendível por cabeça. A intuição é que um mapa captura o sinal principal e o outro captura parte do ruído correlacionado; a subtração tende a reduzir correspondências espúrias. Em faces, esse tipo de mecanismo é interessante porque áreas sem valor genealógico, como fundo da imagem ou cabelo, bem como artefatos espaciais causados por variações de iluminação, podem gerar similaridade alta sem indicar parentesco. No DINOv2-LoRA, essa ideia aparece no módulo de atenção cruzada diferencial entre as duas faces.

\section{Adaptação de baixo posto com LoRA}

Low-Rank Adaptation, ou LoRA \cite{hu2021lora},
é uma técnica de fine-tuning eficiente em parâmetros. Em vez de atualizar toda uma matriz de pesos $W \in \mathbb{R}^{d \times d}$, LoRA preserva $W$ congelada e adiciona um caminho treinável de baixo posto:

\begin{equation}
y = Wx + \frac{\alpha}{r} \, B A x,
\label{eq:lora}
\end{equation}

onde $A \in \mathbb{R}^{r \times d}$ e $B \in \mathbb{R}^{d \times r}$ são matrizes treináveis, $r$ é o posto, com $r \ll d$, e $\alpha$ é um fator de escala. Como apenas $A$ e $B$ recebem gradiente, o número de parâmetros treináveis fica muito menor do que no ajuste fino completo.

Essa técnica é útil quando o backbone é grande e o conjunto de treinamento específico é limitado. Em vez de reescrever todos os pesos do modelo, a adaptação aprende correções de baixa dimensão sobre representações já existentes. No DINOv2-LoRA deste trabalho, LoRA é usado nas projeções qkv e proj dos blocos Transformer, com $r=8$ e $\alpha=16$ na configuração padrão, e funciona como uma forma de adaptar o DINOv2 à tarefa de parentesco facial mantendo o custo de memória mais baixo do que no descongelamento total.

\section{Modelos aumentados por recuperação}

Modelos aumentados por recuperação combinam uma rede treinável com uma memória externa de exemplos. A ideia é que parte da informação útil para uma decisão pode estar em exemplos parecidos já observados, e não apenas nos pesos do modelo. Em vez de decidir somente a partir do par de consulta, o sistema recupera vizinhos em uma galeria e usa esses vizinhos como contexto.

No caso da verificação de parentesco, o uso de recuperação parte do par de consulta $(x_a, x_b)$. Primeiro calcula-se uma assinatura $s_q$ para representar esse par no espaço de embeddings. Depois, em uma galeria pré-computada $\mathcal{G} = \{(s_i, r_i)\}$, em que $s_i$ representa um par positivo do treino e $r_i$ guarda o tipo da relação, são buscados os $K$ pares mais próximos por cosseno. Esses pares recuperados entram como contexto: o módulo de atenção cruzada compara a consulta com os suportes e a saída alimenta tanto a cabeça binária de verificação quanto a cabeça auxiliar de relação.

O interesse desse tipo de modelo está em tornar explícita a comparação com exemplos conhecidos de parentesco. Em uma tarefa com relações variadas, como pais e filhos, irmãos e avós e netos, a memória externa pode oferecer ao modelo casos próximos do par analisado. No Modelo com Recuperação deste trabalho, essa galeria é formada por cerca de 33 mil pares positivos da partição de treino do FIW. Ao mesmo tempo, o método passa a depender da qualidade da galeria, da forma como a assinatura do par é construída e da capacidade de recuperar suportes realmente informativos.

\section{Funções de perda}

A função de perda define o sinal usado para ajustar os parâmetros durante o treinamento. Em verificação binária, a escolha mais direta é a entropia cruzada binária, ou BCE, aplicada ao logit produzido pela cabeça final. Durante a otimização, a minimização da BCE penaliza erros de classificação e fornece gradientes para ajustar os pesos de modo a separar pares positivos e negativos, mas não impõe sozinha uma geometria específica ao espaço de embeddings. Por isso, ela costuma ser combinada ou comparada com perdas de aprendizado métrico.

A Contrastive Loss clássica, proposta por \citeonline{hadsell2006contrastive}, trabalha diretamente sobre pares. A perda aproxima embeddings de pares positivos e afasta pares negativos até uma margem $m$. Em uma versão com similaridade de cosseno, pode ser escrita como

\begin{equation}
\mathcal{L}_{\text{con}} = y \cdot (1 - \cos(z_a, z_b)) + (1 - y) \cdot \max\!\big(0, \cos(z_a, z_b) - m\big),
\label{eq:contrastive}
\end{equation}

Nessa expressão, $y=1$ indica parente e $y=0$ indica não-parente. Os vetores $z_a$ e $z_b$ são embeddings normalizados em $L^2$, e $m$ define a margem a partir da qual o par negativo deixa de contribuir para a perda.

A Triplet Loss \citeonline{schroff2015facenet} usa trios em vez de pares. Cada trio contém uma âncora $z_a$, um positivo $z_p$ da mesma classe e um negativo $z_n$ de classe diferente. A perda força a distância entre âncora e positivo a ser menor que a distância entre âncora e negativo por uma margem mínima. Em reconhecimento facial, essa formulação ficou conhecida por organizar o espaço de embeddings a partir de comparações relativas.

A Supervised Contrastive Loss \cite{khosla2020supervised} generaliza a ideia contrastiva para múltiplos positivos dentro do mesmo batch. Na formulação canônica, a perda usa similaridade cosseno escalada por temperatura e soma a contribuição de todos os positivos da amostra, sem margem explícita. No código experimental deste trabalho, porém, as execuções inicialmente registradas como supervisionadas contrastivas usam distância de cosseno entre pares e margem $m = 0{,}3$. Por equivalência operacional, essa perda é referida nos capítulos seguintes como Contrastive Loss com cosseno e margem, reservando o nome Supervised Contrastive para a formulação canônica de Khosla.

Essas perdas não mudam apenas o valor do erro no treinamento. Elas influenciam a forma do espaço aprendido. Perdas com margem tendem a criar separação mais explícita entre positivos e negativos, enquanto a BCE atua diretamente na decisão de classe. No problema de parentesco facial, essa diferença importa porque o modelo precisa aprender tanto a classificar pares quanto a ordenar corretamente os escores de similaridade.

\section{Modelos multimodais visão-linguagem}

Modelos multimodais visão-linguagem, ou VLMs, combinam representação visual com linguagem natural. CLIP \cite{radford2021learning} foi um dos primeiros modelos amplamente usados nesse formato, treinando um encoder visual e um encoder textual para aproximar imagens e textos correspondentes no mesmo espaço. Modelos mais recentes, como GPT-4V e Claude Sonnet, ampliam esse regime para diálogo multimodal, em que o usuário envia uma imagem junto com uma instrução e recebe uma resposta em linguagem natural.

No parentesco facial, os VLMs entram pela possibilidade de testar a tarefa sem treinamento específico no FIW. O procedimento é simples: envia-se uma imagem com duas faces e um prompt que pergunta se as pessoas são biologicamente aparentadas. A resposta do modelo passa a depender das informações extraídas da imagem e da instrução textual. Essa simplicidade, porém, também limita a leitura do resultado. Como o modelo não foi ajustado para parentesco facial, seu desempenho pode ser influenciado por variáveis de confusão mais imediatas, como idade aparente, gênero, qualidade da imagem, iluminação ou semelhança geral entre as faces.

Nas linhas de base finais deste trabalho, os VLMs são avaliados em verificação binária zero-shot. O prompt exige uma resposta estruturada, com uma decisão kin/non-kin, ou equivalente booleana, e uma confiança no intervalo $[0,1]$. Essa confiança pode ser combinada com a decisão para formar um escore contínuo: respostas positivas recebem escores acima de 0,5 e respostas negativas recebem escores abaixo de 0,5. Assim é possível calcular métricas como ROC AUC, Average Precision e TAR@FAR, ainda que o modelo não tenha um conjunto de validação próprio para ajustar limiar.

A avaliação multiclasse anterior, em que o VLM escolhia uma entre onze relações positivas, permanece apenas como análise diagnóstica. Ela ajuda a observar confusões entre pai-filho, irmãos e avós-netos, mas não substitui a verificação binária porque não apresenta pares non-kin ao modelo.

\section{Métricas de avaliação}

A avaliação dos modelos supervisionados parte da matriz de confusão da tarefa binária. Um verdadeiro positivo ocorre quando um par de parentes é aceito como parente. Um falso positivo ocorre quando um par sem parentesco é aceito por engano. Um verdadeiro negativo é um par sem parentesco corretamente rejeitado, e um falso negativo é um par de parentes rejeitado pelo modelo. A partir desses quatro casos são calculadas as métricas usadas no Capítulo~\ref{cap:Metodologia}.

A acurácia mede a fração total de decisões corretas. Embora seja uma métrica de fácil interpretação, ela pode mascarar vieses de desempenho ou erros concentrados em classes mais difíceis. A precisão mede, entre os pares classificados como parentes, quantos realmente eram positivos. A revocação mede, entre os pares que eram parentes, quantos foram recuperados pelo modelo. O F1 combina precisão e revocação por média harmônica, sendo útil quando se quer observar os dois efeitos em uma única medida. Neste trabalho, essas métricas dependem do limiar $\tau^*$ escolhido na validação.

A métrica principal é o ROC AUC. A curva ROC mostra como a taxa de verdadeiros positivos muda em relação à taxa de falsos positivos quando o limiar varia. O AUC resume essa curva em um único número e pode ser lido como a capacidade do modelo de dar escore maior a um par positivo do que a um par negativo. Por não depender de um limiar específico, ele é adequado para comparar modelos que produzem escores em escalas diferentes.

Average Precision também é independente de um único limiar, mas considera a relação entre precisão e revocação. Ela resume a curva precisão-revocação e é útil quando interessa saber se os pares positivos aparecem bem ranqueados entre os maiores escores. No contexto deste trabalho, ela complementa o AUC porque avalia a qualidade do ranking por outro ponto de vista.

Como a tarefa se aproxima de um problema biométrico de verificação, também são usadas FAR, TAR e TAR@FAR. FAR é a taxa de falsa aceitação, isto é, a proporção de pares negativos aceitos como positivos. TAR é a taxa de aceitação verdadeira, isto é, a proporção de pares positivos aceitos corretamente. A métrica TAR@FAR fixa primeiro uma taxa máxima de falsa aceitação e mede qual TAR o modelo consegue manter nesse ponto. Este trabalho usa FAR igual a $0{,}1$, $0{,}01$ e $0{,}001$. Quanto menor o FAR, mais rígido é o regime de decisão, o que é importante em contextos nos quais falso positivo tem custo alto.

Além das métricas globais, o trabalho reporta acurácia por tipo de relação. Esse recorte mostra se o modelo funciona de forma parecida para pai-filho, mãe-filha, irmãos, avós e netos, ou se o bom resultado geral está concentrado apenas em algumas relações. Esse ponto é importante no FIW porque as classes de avós e netos têm menos pares e tendem a ser mais instáveis.

Para os VLMs binários, as mesmas métricas de verificação podem ser calculadas a partir do escore derivado da decisão e da confiança. A diferença é que não há treino nem validação para escolher um limiar ótimo; por isso, acurácia, precisão, revocação e F1 usam o limiar natural de 0,5, enquanto AUC, Average Precision e TAR@FAR são lidos como medidas mais comparáveis entre modelos. Para a avaliação multiclasse histórica dos VLMs, quando mencionada, continuam valendo acurácia global, macro-precision, macro-recall, macro-F1 e matriz de confusão entre as onze relações positivas.

Com essas definições, a leitura dos resultados fica menos dependente de uma única medida. O AUC e a Average Precision avaliam a ordenação geral dos pares, as métricas com limiar mostram o comportamento no ponto escolhido, e o TAR@FAR mostra o desempenho quando a falsa aceitação é controlada explicitamente. Nos VLMs, a ausência de calibração supervisionada e a distribuição discreta de confiança são parte da interpretação, porque podem limitar a operação em regimes de baixa falsa aceitação.
